\citet{GentzkowShapiro:2014} give language-agnostic rules for research code.
This entry lists how the template implements them in R, Python and Julia.

\begin{description}
	\item[Automation.]
	One command rebuilds everything: the top-level \texttt{make} runs every task, then the paper, slides and logbook.
	There are no manual export, copy-paste or point-and-click steps.
	Raw formats are converted in code (\texttt{readxl}, \texttt{pandas.read\_excel}, \texttt{XLSX.jl}).
	Tables and in-text numbers are written by code straight to \texttt{.tex} files and never copied from a console.

	\item[Version control.]
	Code and documents are tracked with git and GitHub.
	Before committing, run \texttt{make clean} and \texttt{make} in every changed task and everything downstream.
	Unlike Gentzkow and Shapiro's setup, data are not committed: raw data are referenced read-only, and outputs are rebuilt.

	\item[Directories.]
	Each task separates \texttt{input/}, \texttt{code/} and \texttt{output/}.
	Inputs are symbolic links and paths are relative, so the project is portable.
	\texttt{make clean} clears outputs so that a rebuild proves they come from the current code.

	\item[Keys.]
	Every saved table has a unique, non-missing key, checked by \texttt{assert\_key()} inside \texttt{save\_table()}.
	Data stay normalised as long as possible, built in three steps: normalised tables that preserve the source; the same tables with derived variables added; and a final merge into the analysis array with no data manipulation.
	Each variable lives in the table for the level it describes.

	\item[Abstraction.]
	Write a function the second time an operation recurs, not before.
	Functions used across tasks go in \texttt{lib/} with unit tests.

	\item[Documentation.]
	Code should explain itself.
	Comments explain why and cite sources.
	Constants are computed from their named components in code rather than restated in comments.
	Guards are code (\texttt{stopifnot}, \texttt{assert}, \texttt{@assert}), not ``don't do this'' comments.

	\item[Management.]
	Tasks are managed in GitHub Issues: one issue per task, one branch per issue, and a pull request that closes the issue.

	\item[Code style.]
	Lines are at most 100 characters (checked by \texttt{lintr}, \texttt{ruff} and JuliaFormatter).
	Functions are short ($\lesssim 80$ lines), take explicit inputs and return explicit outputs, and are ordered for top-to-bottom reading.
	Names are descriptive.
	Logical tests say what they mean.
	Error checks target costly failures and silent wrong answers, and repeated validation is centralised in one function.
	Unit tests replace manual checks at the console.
	Slow code is profiled.
	Slow tasks (estimation, simulation) are separate from fast ones (tables, figures) and store everything they compute.
\end{description}

Dropped as Stata-specific or obsolete: Stata examples, Stat/Transfer, \texttt{outreg}, Windows \texttt{.bat} run scripts (replaced by Make), and the tool recommendations (Subversion, JIRA, Asana, Evernote), which git, GitHub Issues and this logbook replace.
